\documentclass[11pt]{article}
\usepackage{amsmath, amssymb, booktabs, hyperref}
\usepackage[margin=1in]{geometry}

\title{Mixed-Effects Analysis of Multi-Trial Forecasting Results\\
on the AIBQ2 Benchmark}
\author{Kevin Murphy (with help from Claude Code)}
\date{6 April 2026}

\begin{document}
\maketitle

\section{Introduction}

We evaluate an LLM-based agentic forecaster on the AIBQ2 benchmark
(113 binary questions from the Q2~2025 Metaculus AI Benchmark Tournament).
Each base model is run $K=5$ independent trials per question, and the
trials can be aggregated in different ways.

We conduct two analyses:
\begin{itemize}
  \item \textbf{Simple analysis} (Section~\ref{sec:simple}):
    4 methods varying thinking level (default vs.\ high) and calibration
    (on/off), all using the same shrinkage aggregation over $K=5$ trials.
    This isolates the effect of thinking and calibration.
  \item \textbf{Expanded analysis} (Section~\ref{sec:expanded}):
    8 methods that also vary the aggregation scheme
    (\texttt{mean:1}, \texttt{mean:5}, \texttt{shrink:5}),
    giving a factorial design that quantifies the contribution of each factor.
\end{itemize}


\section{Model}

We fit a two-way fixed-effects model to the per-question, per-method scores:
\begin{equation}
  \mathrm{score}(q, m, k) = \mu + \alpha_m + \gamma_q + \varepsilon_{q,m,k},
  \label{eq:model}
\end{equation}
where:
\begin{itemize}
  \item $\mu$ is the grand mean score across all methods and questions,
  \item $\alpha_m$ is the \emph{method effect} for method $m$ (what we want to estimate),
  \item $\gamma_q$ is the \emph{question difficulty} for question $q$ (nuisance parameter),
  \item $\varepsilon_{q,m,k}$ is the residual trial noise.
\end{itemize}

We use two views of the data:

\paragraph{Aggregated scores (for method effects).}
For each (question~$q$, method~$m$) pair, we have one aggregated score:
\begin{equation}
  s_{q,m} = \ell\!\bigl(\hat{p}_{q,m},\; o_q\bigr),
\end{equation}
where $\hat{p}_{q,m}$ is the aggregated forecast, $o_q \in \{0,1\}$ is the
outcome, and $\ell$ is the scoring function (Metaculus or Brier).

For methods that aggregate $K$ trials (e.g.\ \texttt{mean:5}, \texttt{shrink:5}),
the score is computed from the aggregated forecast:
\begin{equation}
  s_{q,m} = \ell\!\Bigl(\mathrm{agg}\bigl(p_{q,m,1}, \ldots, p_{q,m,K}\bigr),\; o_q\Bigr).
\end{equation}

For \texttt{mean:1}, there is no aggregation, so we report the expected
single-trial score, averaged over all $K$ individual trials:
\begin{equation}
  s_{q,m} = \frac{1}{K} \sum_{k=1}^{K} \ell(p_{q,m,k},\; o_q).
  \label{eq:mean1}
\end{equation}
Note the distinction: \texttt{mean:K} scores the aggregated forecast
$\ell(\bar{p}, o)$, while \texttt{mean:1} averages the individual scores
$\bar{\ell}(p_k, o)$.  These differ because $\ell$ is nonlinear.

We fit model~\eqref{eq:model} to the $Q \times M$ aggregated scores
via alternating projections (block coordinate descent), equivalent to
$\texttt{feols}(\text{score} \sim 1 \mid \text{question} + \text{method})$
in \texttt{pyfixest}:
\begin{align}
  \gamma_q &\leftarrow \frac{1}{M} \sum_m (s_{q,m} - \mu - \alpha_m), \\
  \alpha_m &\leftarrow \frac{1}{Q} \sum_q (s_{q,m} - \mu - \gamma_q),
\end{align}
until convergence (typically $< 50$ iterations).

\paragraph{Trial-level scores (for ANOVA).}
To separate within-trial noise from between-question difficulty,
we use all $K \times M \times Q$ individual trial scores:
\begin{equation}
  s_{q,m,k} = \ell(p_{q,m,k},\; o_q), \quad k = 1, \ldots, K,
\end{equation}
where $p_{q,m,k}$ is the raw forecast from trial~$k$ (before aggregation).
We perform a Type~III sum-of-squares decomposition:
\begin{align}
  \mathrm{SS}_{\mathrm{total}} &= \sum_{q,m,k} (s_{q,m,k} - \bar{s})^2, \\
  \mathrm{SS}_{\mathrm{method}} &= KQ \sum_m (\bar{s}_{\cdot m \cdot} - \bar{s})^2, \\
  \mathrm{SS}_{\mathrm{question}} &= KM \sum_q (\bar{s}_{q \cdot \cdot} - \bar{s})^2, \\
  \mathrm{SS}_{\mathrm{residual}} &= \mathrm{SS}_{\mathrm{total}}
    - \mathrm{SS}_{\mathrm{method}} - \mathrm{SS}_{\mathrm{question}}.
\end{align}


\section{Simple Analysis: 4 Methods}
\label{sec:simple}

The four methods are:
(1)~\textbf{default}: Gemini~3.1~Pro, default reasoning;
(2)~\textbf{high}: high reasoning;
(3)~\textbf{default+cal}: method~1 with Platt calibration (LOO CV);
(4)~\textbf{high+cal}: method~2 with Platt calibration.
All use $K=5$ trials with logit-space shrinkage aggregation.

\subsection{Method effects}

The grand mean across all four methods is $\hat{\mu} = 36.6$.

\begin{table}[h]
\centering
\begin{tabular}{lrr}
\toprule
Method & Effect ($\hat{\alpha}_m$) & Adjusted score \\
\midrule
high + cal & $+3.0$ & $39.6$ \\
high       & $+0.7$ & $37.3$ \\
default + cal & $+0.2$ & $36.8$ \\
default    & $-3.9$ & $32.7$ \\
\bottomrule
\end{tabular}
\caption{Method effects for Metaculus score (higher = better).
  The adjusted score is $\hat{\mu} + \hat{\alpha}_m$.}
\label{tab:simple_effects}
\end{table}

The ranking is: high+cal $>$ high $>$ default+cal $>$ default.
The total spread is $39.6 - 32.7 = 6.9$ points.

\subsection{Variance decomposition}

\begin{table}[h]
\centering
\begin{tabular}{lrrrr}
\toprule
Source & SS & \% & df & $F$ \\
\midrule
Method & $11{,}045$ & $0.0\%$ & 3 & 1.28 \\
Question & $24{,}605{,}626$ & $79.9\%$ & 112 & $76.4^*$ \\
Residual & $6{,}165{,}926$ & $20.0\%$ & $2{,}144$ & --- \\
\midrule
Total & $30{,}782{,}596$ & $100\%$ & $2{,}259$ & \\
\bottomrule
\end{tabular}
\caption{ANOVA for 4 methods ($K=5$, $Q=113$).
  $^*$Significant at $p < 0.05$; $F_{\mathrm{crit}}(3, 2144) \approx 2.6$.}
\label{tab:simple_anova}
\end{table}

The key finding: \textbf{question difficulty explains $80\%$ of total
variance} (Table~\ref{tab:simple_anova}), trial-to-trial stochasticity
explains $20\%$, and method differences are negligible
($< 0.1\%$, $F = 1.28$, not significant at $p = 0.05$).

The question effect is highly significant
($F = 76.4 \gg F_{\mathrm{crit}}$), confirming that questions genuinely
differ in difficulty.
\begin{itemize}
  \item Whether a question is ``easy'' or ``hard'' matters far more than
    which method you use.
  \item Within-trial noise ($20\%$) is substantial but secondary to question
    difficulty---the multi-trial aggregation ($K=5$) successfully reduces
    this component.
\end{itemize}

\subsection{Pairwise comparisons}

We compare each method to the best (high+cal) using paired differences
on the same 113 questions, with bootstrap $p$-values.
Specifically, for each question $q$ we compute
$\Delta_q = s_{q,m} - s_{q,\mathrm{ref}}$, then bootstrap the mean
of $\Delta_q$ over 2{,}000 resamples.  The quantity $p(\Delta > 0)$ is the
fraction of bootstrap replicates where the mean difference is positive
(i.e., the method beats the reference).

\begin{table}[h]
\centering
\begin{tabular}{lrrr}
\toprule
Method vs.\ high+cal & $\Delta$ mean & $\Delta$ std & $p(\Delta > 0)$ \\
\midrule
default           & $-6.9$ & $46.8$ & $0.066$ \\
default + cal     & $-2.8$ & $23.3$ & $0.104$ \\
high              & $-2.3$ & $41.9$ & $0.283$ \\
\bottomrule
\end{tabular}
\caption{Pairwise comparisons vs.\ high+cal (Metaculus score, $n=113$).
  $p(\Delta > 0)$ is the bootstrap probability that the method beats
  the reference.}
\label{tab:simple_pairwise}
\end{table}

No pairwise comparison reaches $p < 0.05$
(Table~\ref{tab:simple_pairwise}).  The largest effect
(default vs.\ high+cal, $\Delta = -6.9$) has $p = 0.066$, driven by the
high per-question standard deviation ($\sigma = 46.8$).  With 113 questions,
we are underpowered to detect effects smaller than $\sim 9$ points at
$\alpha = 0.05$.

\subsection{Brier score}

\begin{table}[h]
\centering
\begin{tabular}{lrr}
\toprule
Method & Effect & Adjusted \\
\midrule
high + cal & $-0.0064$ & $0.1373$ \\
high       & $-0.0046$ & $0.1391$ \\
default + cal & $+0.0008$ & $0.1446$ \\
default    & $+0.0101$ & $0.1538$ \\
\bottomrule
\end{tabular}
\caption{Brier score method effects (lower = better).}
\label{tab:brier_effects}
\end{table}

The Brier score results (Table~\ref{tab:brier_effects}) mirror the
Metaculus ranking.  The ANOVA is similar: question difficulty $\sim 80\%$,
trial noise $\sim 20\%$, method $< 0.1\%$ (not significant).


\section{Expanded Analysis: 8 Methods}
\label{sec:expanded}

We now include aggregation variants, giving 8 methods:
2 thinking levels $\times$ 3 aggregation schemes (\texttt{mean:1},
\texttt{mean:5}, \texttt{shrink:5}), plus 2 calibrated variants
(\texttt{shrink:5+cal}).

For \texttt{mean:1}, the score is the expected single-trial score
$\mathbb{E}_k[\ell(p_{q,m,k}, o_q)]$, averaged over all $K$ trials.
This represents the performance of a single run without any aggregation.

\subsection{Method effects and ANOVA}

\begin{table}[h]
\centering
\begin{tabular}{lrr}
\toprule
Method & Effect ($\hat{\alpha}_m$) & Adjusted score \\
\midrule
high / shrink:5 + cal & $+6.7$ & $39.6$ \\
high / shrink:5 & $+4.4$ & $37.3$ \\
default / shrink:5 + cal & $+3.9$ & $36.8$ \\
high / mean:5 & $+2.5$ & $35.4$ \\
default / shrink:5 & $-0.2$ & $32.7$ \\
default / mean:5 & $-2.3$ & $30.6$ \\
high / mean:1 & $-5.3$ & $27.6$ \\
default / mean:1 & $-9.7$ & $23.2$ \\
\bottomrule
\end{tabular}
\caption{Method effects for all 8 variants (grand mean $= 32.9$).}
\label{tab:expanded_effects}
\end{table}

\begin{table}[h]
\centering
\begin{tabular}{lrrrr}
\toprule
Source & SS & \% & df & $F$ \\
\midrule
Method & $24{,}282$ & $0.3\%$ & 7 & $3.93^*$ \\
Question & $7{,}069{,}749$ & $90.8\%$ & 112 & $71.56^*$ \\
Residual & $691{,}561$ & $8.9\%$ & $784$ & --- \\
\midrule
Total & $7{,}785{,}592$ & $100\%$ & $903$ & \\
\bottomrule
\end{tabular}
\caption{ANOVA for 8 methods ($^*$significant at $p < 0.05$;
  $F_{\mathrm{crit}}(7, 784) \approx 2.0$).}
\label{tab:expanded_anova}
\end{table}

With 8 methods, the method effect is now \textbf{significant} ($F = 3.93$,
$p < 0.05$), driven by the $16.4$-point spread between \texttt{mean:1}
and \texttt{shrink:5+cal}.  The residual dropped from $20\%$ (4 methods)
to $8.9\%$ because the aggregation variants absorb what was previously
attributed to trial noise.

\subsection{Factorial decomposition}

We decompose the method effects into three factors by averaging each
factor's effect over the levels of the other factors:

\begin{enumerate}
  \item \textbf{Aggregation} (mean:1 $\to$ mean:5 $\to$ shrink:5),
    averaged over thinking levels:
    \begin{itemize}
      \item mean:1 $\to$ mean:5:
        $\frac{(30.6 - 23.2) + (35.4 - 27.6)}{2} = +7.6$
      \item mean:5 $\to$ shrink:5:
        $\frac{(32.7 - 30.6) + (37.3 - 35.4)}{2} = +2.0$
      \item Total: $+9.6$ points
    \end{itemize}
  \item \textbf{Thinking} (default $\to$ high),
    averaged over aggregation methods:
    \[
      \frac{(27.6 - 23.2) + (35.4 - 30.6) + (37.3 - 32.7)}{3} = +4.6
    \]
  \item \textbf{Calibration} (shrink:5 $\to$ shrink:5+cal),
    averaged over thinking levels:
    \[
      \frac{(36.8 - 32.7) + (39.6 - 37.3)}{2} = +3.2
    \]
\end{enumerate}

\begin{table}[h]
\centering
\begin{tabular}{lr}
\toprule
Factor & Gain (Metaculus points) \\
\midrule
Aggregation (mean:1 $\to$ shrink:5) & $+9.6$ \\
\quad of which mean:1 $\to$ mean:5 & $+7.6$ \\
\quad of which mean:5 $\to$ shrink:5 & $+2.0$ \\
Thinking (default $\to$ high) & $+4.6$ \\
Calibration (shrink:5 $\to$ shrink:5+cal) & $+3.2$ \\
\midrule
Sum of individual factors & $+17.4$ \\
Observed total (default/mean:1 $\to$ high/shrink:5+cal) & $+16.4$ \\
\bottomrule
\end{tabular}
\caption{Incremental gains from each factor.  The sum ($17.4$) slightly
  exceeds the observed total ($16.4$) due to a small negative interaction
  between factors.}
\label{tab:factorial}
\end{table}


\section{Discussion}

\paragraph{Ranking of interventions.}
The three factors ranked by Metaculus score improvement:
\begin{enumerate}
  \item Multi-trial aggregation: $+9.6$ points ($5\times$ compute cost)
  \item High thinking: $+4.6$ points (no extra cost per trial)
  \item Calibration: $+3.2$ points (negligible cost)
\end{enumerate}

\paragraph{Statistical significance.}
With 4 methods (Section~\ref{sec:simple}), the method $F$-ratio is $1.28$
(not significant): the methods are too similar.
With 8 methods (Section~\ref{sec:expanded}), $F = 3.93$ (significant):
the aggregation variants widen the spread enough to detect.  Question
difficulty dominates in both analyses ($80$--$91\%$ of variance).

\paragraph{Practical implications.}
Calibration and high thinking are ``free'' improvements ($+3.2$ and $+4.6$
points respectively) that should always be used.  Multi-trial aggregation
is the largest single intervention ($+9.6$), but costs $5\times$ in compute.
For a fixed budget, the optimal allocation is: use high thinking (free),
calibrate (free), then spend remaining budget on as many trials as
affordable.

\end{document}
